\documentclass[10pt]{letter}
\usepackage[margin=2.5cm, top=2cm, bottom=3cm]{geometry}
\usepackage{xspace}
\usepackage[hidelinks]{hyperref}

\newcommand{\HW}{\textsc{HiveWatch}\xspace}
\newcommand{\titleofourpaper}{[Full Manuscript Title]}

\pagestyle{empty}
\setlength{\parindent}{0pt}
\setlength{\parskip}{0.7em}

\begin{document}

September 22, 2026

Editor-in-Chief\\
SoftwareX

Dear Editor-in-Chief,

We are submitting our manuscript titled
``\textbf{\texttt{HiveWatch}: Framework-Agnostic Observability for Federated Learning}'' for consideration for publication
in SoftwareX. We confirm that this work is original, is not
under consideration elsewhere, and has been approved for
submission by all authors.

Federated learning experiments involve distributed clients
and different software frameworks, making consistent monitoring
and experiment tracking challenging. We present \texttt{HiveWatch}, a
lightweight Python package for monitoring federated learning
experiments, with optional integrations for Weights \& Biases
and MLflow. Its core package has no mandatory third-party
dependencies, allowing users to install additional backends
according to their needs.

\texttt{HiveWatch} includes runnable integration examples for APPFL, Flower,
and NVIDIA FLARE, together with developer documentation to
support adoption. The software is available through PyPI and
GitHub under the MIT License. The manuscript identifies
version v0.3.0 as the accompanying release, enabling readers
to access the specific version described in the paper.

We believe that \texttt{HiveWatch's} practical utility, open-source
availability, and support for multiple federated learning
frameworks make it a suitable contribution to SoftwareX.
The software and accompanying examples provide researchers
and practitioners with a reusable tool for monitoring
distributed machine learning experiments.

Thank you for considering our manuscript. We look forward
to your feedback.

Best regards,\\
Abhijit Chunduru\\
On behalf of all authors

\medskip
\underline{Corresponding Author Contact Information}:

\begin{minipage}[t]{0.7\textwidth}
	Ravi Madduri, \textit{Senior Scientist and Group Leader} (madduri@anl.gov) \\
    
	Data Science and Learning Divison,\\
	Argonne National Laboratory\\
    Lemont, IL 60439, United States. \\
\end{minipage}

\end{document}